% Anexos de terceiros utilizados como material complementar.

\anexo[ane:exemplo]{Exemplo de Anexo}

Os anexos são documentos que não são de autoria do autor do TCC, mas que são pertinentes à pesquisa. Eles podem incluir gráficos, documentos oficiais, normativas e textos que foram utilizados como referência ou que fornecem suporte à pesquisa. Segundo a Associação Brasileira de Normas Técnicas (ABNT), ``os anexos são documentos que, embora não sejam parte integrante do texto, são relevantes para a compreensão do tema tratado''~\cite{abnt2012}. O uso de anexos é indicado quando se deseja fornecer ao leitor informações complementares que não foram geradas pelo autor, mas que ajudam a esclarecer ou contextualizar a pesquisa.

Ambos, apêndices e anexos, devem ser listados em uma página específica após a conclusão do trabalho, de forma a facilitar a consulta por parte do leitor. A correta distinção entre apêndices e anexos, bem como o seu uso apropriado, é essencial para a organização e a credibilidade do TCC. Ao fornecer informações adicionais de maneira clara e sistemática, o autor contribui para uma melhor compreensão do seu trabalho, enriquecendo a experiência do leitor.
